\section{Introdução}
\label{sec:introducao}

O trabalho segue a Variante A (ambiente interno), pois a soma dos registros acadêmicos (RAs), 977818, termina em dígito par. Reconstruímos metricamente uma cozinha a partir de um vídeo monocular, da seleção de quadros à nuvem esparsa, às nuvens densas, à planta baixa e à malha~\cite{hartley2004multiple,schoenberger2016sfm}. O vídeo e os quadros usados estão descritos na Seção~\ref{sec:aquisicao}. % fonte: notas/t1_metadados.md; cozinha/config.resolved.yaml; sfm_metrics.csv

O método (Seção~\ref{sec:metodo}) usa SIFT (Scale-Invariant Feature Transform) e estrutura a partir do movimento (SfM, Structure from Motion) com o mapeador global GLOMAP~\cite{pan2024glomap}.
Os resultados estão na Seção~\ref{sec:resultados}. O código está no repositório \href{https://github.com/oliveiraMats2/3d_indoor_environment}{\mychgblue{\texttt{3d\_indoor\_environment}}}. Todas as figuras em alta resolução estão na pasta \href{https://drive.google.com/drive/folders/1zD92kOM3-jJZeM48MSD7dSlieicfB6Qa}{\mychgblue{Image 3D Report}}; clicar em cada figura abre a sua versão no Drive. O vídeo, as nuvens, a malha, as medidas, as métricas, a ablação, a planta e a comparação com o LiDAR exigidos pelo enunciado estão na subpasta \href{https://drive.google.com/drive/folders/1tiqDEpNtPre5AfVaHxmzFQ1o4eBh0snh}{\mychgblue{Requisitos T2}}.
